\section{Cái chương này không đo được, và vì sao tôi không đo bừa}
\label{sec:ch09-khong-do-duoc}

\index{tiền huấn luyện!corpus của sách không với tới}%
Mọi chương từ chương~\ref{ch:encoder-decoder} tới đây đều kết thúc bằng một
bảng tôi tự chạy. Chương này thì không, và tôi nợ bạn lý do: \enquote{không đo}
ở đây là một quyết định đã cân nhắc, và mục này là chỗ trình bày nó.

Corpus của sách sinh từ một văn phạm phi ngữ cảnh nhỏ. Nó có ba tính chất mà
\code{toy\_corpus.py} ghi trong docstring của mình từ chương~\ref{ch:encoder-decoder}:
văn phạm hữu hạn nên một mô hình học được nó chính xác, từ điển đóng nên không
có từ lạ, và câu thì vô nghĩa về mặt ngữ nghĩa.

Ba tính chất ấy đã lấy đi của sách hai thứ. Chương~\ref{ch:mot-vector-cho-moi-tu}
không làm phép cắt bỏ hai chiều nói được gì, vì văn phạm không có nhập nhằng nào
để một cơ chế khử nhập nhằng thể hiện. Chương~\ref{ch:transformer} thì gặp giới
hạn ở chỗ khác và nhẹ hơn.

Ở chương này nó lấy đi cả thí nghiệm.

\subsection*{Vì sao một thí nghiệm tiền huấn luyện trên corpus ấy đo số không}

\index{tiền huấn luyện!vì sao không có thí nghiệm}%
\tn{Tiền huấn luyện}{pretraining} được cái gì đó chỉ khi pha một thấy thứ mà
pha hai cần và không có đủ. Nói cho gọn: hai pha phải nhìn hai phân phối khác nhau, và cái học được
ở phân phối rộng phải chuyển sang được phân phối hẹp.

Tách corpus của sách làm hai phần và gọi phần lớn là \enquote{tiền huấn luyện}
thì hai phần ấy rút từ đúng một văn phạm, với đúng một từ điển, và đúng một tập
luật. Không có phân phối thứ hai. Một mô hình chạy pha một rồi pha hai trên đó
chỉ đang huấn luyện lâu hơn, và bảng kết quả sẽ cho thấy nó tốt hơn mô hình
huấn luyện ngắn, đúng như mọi bảng epoch trong sách này đã cho thấy từ
chương~\ref{ch:transformer}.

Con số sẽ có. Nó sẽ đo số bước, và mục~\ref{sec:ch07-corpus} đã đo số bước rồi.
Cái nó không đo là chuyển giao, vì không có gì để chuyển.

Sửa cho đúng thì phải đổi corpus, không phải đổi script. Cần một văn phạm sinh
ra một miền rộng và một miền hẹp nằm trong nó, sao cho miền hẹp có ít dữ liệu
và có cấu trúc mà chỉ miền rộng dạy được. Đó là một corpus khác hẳn corpus sách
này đang dùng, và dựng nó là công việc của cả một phiên làm việc, không phải
của một mục.

Tôi viết chỗ này ra vì cái cám dỗ ngược lại rất mạnh. Một bảng có số trông
giống một chương có bằng chứng, và một chương toàn trích dẫn trông mỏng hơn.
Nhưng một bảng đo sai đại lượng thì tệ hơn không có bảng: nó đưa cho người đọc
một con số để nhớ, và con số ấy nói về chuyện khác.

\subsection*{Chỗ duy nhất sách này chạm được vào luận điểm, và nó ở chương trước}

\index{tiền huấn luyện!dropout đổi chiều theo lượng dữ liệu}%
Có đúng một phép đo trong sách mà kết luận của nó lật chiều chỉ vì lượng dữ
liệu, và nó không nằm ở chương này.

Mục~\ref{sec:ch08-dieu-chuan} chạy dropout ở đúng giá trị bài Vaswani dùng, trên
corpus của sách, và nó làm hỏng. Bài đo cùng cơ chế ấy, cùng giá trị tham số
ấy, và nó đáng 1.2 BLEU. Cùng công thức, hai kết luận ngược nhau, và cái phân
biệt là 6000 câu sinh từ văn phạm so với 4.5 triệu cặp câu thật.

Đó là luận điểm của chương này ở dạng nhỏ nhất mà sách này chạm được: một quyết
định kỹ thuật đúng ở quy mô này và sai ở quy mô kia, không phải vì công thức
đổi mà vì lượng dữ liệu đổi.

Bài BERT và bài GPT là cùng câu chuyện ấy ở quy mô mà nó không còn là một chi
tiết điều chuẩn nữa mà thành cả kiến trúc của quy trình. Và
chương~\ref{ch:vit} là nơi nó đạt dạng gọn nhất: cùng một mô hình, không đổi
một dòng mã, thua ResNet trên ImageNet-1k và thắng sau khi được tiền huấn luyện
trên JFT-300M~\autocite{dosovitskiy2021vit}.

Khoảng cách giữa 6000 câu sinh từ văn phạm và 300 triệu ảnh có nhãn là khoảng
cách mà sách này không bắc qua được bằng một thí nghiệm. Phần ấy là trích dẫn,
và mục này là chỗ nói ra điều đó thay vì để bạn tự đoán.
